\section{Introduction}
\label{sec:introduction}

A fermionic state can be compared with simple states in two different ways.
One can measure its overlaps with them, or ask how many are needed to express
it as a sum. The first question concerns a probability distribution; the
second concerns an exact vector identity. This paper treats both questions
using the geometry of Slater and Gaussian states.

For a fixed number of particles, the one-particle space is $\C^n$ and the
$p$-particle space is its exterior power $\Lambda^p\C^n$. A unit vector
$u_1\wedge\cdots\wedge u_p$, with the $u_i$ orthonormal, is called a
\emph{Slater vector}. It describes the occupation of $p$ mutually orthogonal
one-particle modes. Even in a small space, most vectors need not have this
form. For example, in $\Lambda^2\C^4$ the normalized vector
\begin{equation}\label{intro:two-particle}
 \psi=\frac{e_1\wedge e_2+e_3\wedge e_4}{\sqrt2}
\end{equation}
is not a Slater vector: its exterior square is nonzero, whereas the square
of a decomposable two-vector vanishes. Its overlap with $e_1\wedge e_2$ has
squared modulus $1/2$. Rotating the reference Slater vector gives an entire
distribution of such overlaps.

A density operator is a positive semidefinite operator $\rho$ of trace one;
it includes both pure states $|\psi\rangle\langle\psi|$ and mixtures. If
$\Omega_x$ is a unit coherent vector indexed by a point $x$ of the coherent
orbit, its \emph{Husimi function} is
\[
 Q_\rho(x)=\langle\Omega_x,\rho\Omega_x\rangle.
\]
The inner product is conjugate-linear in its first argument. The orbit
carries the probability measure $\nu$ induced by normalized Haar measure
of the acting compact group. Thus $Q_\rho$ is a random variable in $[0,1]$.
For a $D$-dimensional irreducible representation, its mean is $1/D$.
For a unit input vector $\psi$, we abbreviate
$Q_\psi=Q_{|\psi\rangle\langle\psi|}$.

\subsection{Comparing the whole distribution}
Our distribution theorem says that coherent inputs maximize every continuous
convex average:
\begin{equation}\label{intro:convex-order}
 \int \Phi(Q_\rho)\dd\nu
 \leq \int \Phi(Q_{P_\Omega})\dd\nu,
 \qquad P_\Omega=|\Omega\rangle\langle\Omega|.
\end{equation}
Here $\Omega$ is any coherent unit vector, and $\Phi$ is continuous and convex
on $[0,1]$. This relation is called \emph{convex order}. Since all these
random variables have the same mean, it says that the coherent distribution
has the greatest spread in every convex test. It implies a comparison of
variances by taking $\Phi(t)=t^2$, and of Wehrl entropies by taking
$\Phi(t)=t\log t$ with the continuous value zero at $t=0$. For every
strictly convex $\Phi$, equality in \eqref{intro:convex-order} forces the
whole input density to be a pure coherent state.

We prove this statement for all basic exterior powers
(Theorem~\ref{ext:convex-order}), for both basic half-spin representations
(Theorem~\ref{spin:convex-order}), and for the basic odd-dimensional spin
representation (Corollary~\ref{spin:odd}). In the half-spin realization
used below, the coherent vectors are pure fermionic Gaussian vectors,
whose particle number may vary while their parity remains fixed.
Section~\ref{sec:spin} constructs this realization and explains its
relation to odd-dimensional spin.

The entropy problem has a long history. Gnutzmann and \.{Z}yczkowski
introduced R\'enyi--Wehrl localization measures in the generalized coherent
state setting~\cite{gnutzmann2001}. Sugita proved coherent-state optimality
for integer indices greater than one~\cite{sugita2002}, and studied the
resulting moments for many-particle states~\cite{sugita2003}. Lieb and
Solovej proved a convex-function comparison for symmetric $\mathrm{SU}(N)$
coherent states~\cite{liebsolovej2016}. Symmetric powers describe a different
representation series from the exterior powers considered here. Our
comparison treats arbitrary continuous convex functions and all mixed
inputs in the stated fermionic series. In particular, it yields all real
positive-order R\'enyi--Wehrl minima, together with their equality cases.

The main step is a compatibility statement about conditioning. Detecting one
occupied mode turns the original density into a density on a smaller
fermionic space. Its normalization is an occupation probability, determined
by a one-particle matrix or a Majorana covariance matrix. We express this
matrix as a convex combination of coherent matrices using one set of
probabilities, chosen from the diagonal of the same rotated full density.
These probabilities work for every subsequent measured direction. The
conditional coherent states also remain coherent in the smaller space.
The induction therefore uses the same comparison state at every stage.
The equality argument recovers the full density from this common
decomposition; it does not identify a many-particle state merely by its
one-particle matrix.

\subsection{Counting Gaussian summands}
Allowing superpositions of particle numbers leads to a different question.
The \emph{Gaussian rank} $\chi_G(\psi)$ is the least number of nonzero
Gaussian vectors whose sum is $\psi$, with arbitrary complex coefficients
absorbed into the summands. The definition and parity convention appear in
Section~\ref{sec:gaussian-rank}.

On four modes, let $|0000\rangle$ and $|1111\rangle$ denote the vacuum and
the state with all four modes occupied, and put
\[
 M=\frac{|0000\rangle+|1111\rangle}{\sqrt2}.
\]
This state is a basic non-Gaussian resource in matchgate computation. Its
two-copy expansion has four Gaussian terms. A shorter expansion, however,
could use terms coupling the two blocks of four modes. Cudby and Strelchuk
proved a rank-four result under symmetry restrictions and formulated the
unrestricted two-copy problem~\cite{cudby2023}. Gaussian decompositions are
also central to the simulation methods of Dias and Koenig~\cite{diaskoenig2024}.
Related multiplicativity results for Gaussian fidelity and extent concern
different optimization problems~\cite{reardon2024}.
Recent work of Tarabunga studies non-Gaussianity through the eigenvalue
distribution of a two-copy Majorana operator and its behavior under
postselected Gaussian operations~\cite{tarabunga2026}.

We prove that $\chi_G(M^{\otimes2})=4$, and more generally that Gaussian
rank is multiplicative for two arbitrary nonzero four-mode even vectors.
The obstruction is a polynomial built from the full bilinear Gram of the
Clifford bivector action. It is nonzero at the target and vanishes on every
sum of three pure spinors. A common change of coordinates first normalizes
two summands. A polynomial deformation then brings the third summand into
a tractable chart, and a specialization argument recovers all degenerate
cases. The change of coordinates acts on the entire sum: transforming the
summands independently would destroy the equation being tested.

\subsection{Formal verification and organization}
The accompanying Lean development verifies the exterior-power
convex-order theorem, the actual Slater-orbit probability measure, the
strict equality characterization, and the positive-order moment and entropy
minimizer theorems. These are end-to-end proofs with no additional
mathematical hypotheses beyond those stated in the theorems. The standard
logical foundations and the precise boundary of the formalization are
recorded in Section~\ref{sec:formalization}. In particular, the Gaussian-rank
four conclusion still has a paper proof supplemented by verified coordinate
certificates and original-object lemmas.

The sources, full declaration types, and verification evidence are available at
\begin{center}
 \small\url{https://github.com/veridiscoverLab/fermionic-coherent-lean}.
\end{center}
We begin with exterior powers, where the conditioning argument is most
transparent, and then pass to spin representations. After evaluating the
coherent distributions and entropy constants, we turn to Gaussian rank.
The final application section gives quantitative entropy bounds and the
information capacity of the specified coherent measurements.
